% ============================================================================
% Fichier  : front/01_avant_propos.tex (Avant-propos)
% Version  : 1.0 -- novembre 2026
% ============================================================================

\chapter*{Avant-propos}
\addcontentsline{toc}{chapter}{Avant-propos}
\markboth{Avant-propos}{Avant-propos}

Le 23 juin 2026, le délai de dix-huit mois que la loi n°~2024/017 accordait
aux organisations camerounaises pour se mettre en conformité a expiré. Le
lendemain, une microfinance de Bafoussam reçoit la lettre d'un client dont
le dossier a fuité quinze mois plus tôt. Il demande ce que l'établissement
sait de lui, à qui il l'a transmis, et pourquoi personne ne l'a prévenu. Le
directeur général pose alors à ses équipes la question qui ouvre ce cours~:
sommes-nous en règle~?

Répondre exige un audit. Un audit de protection des données ne se contente
pas de lire des procédures~: il recense les traitements réels, les compare
à ce que la loi permet, vérifie que les personnes peuvent exercer leurs
droits, que les données sont protégées, que les sous-traitants et les
transferts sont encadrés, puis il conclut. Ce cours enseigne cette démarche,
dans le cadre de la loi camerounaise, avec les normes internationales qui
l'outillent.

Il s'adresse aux étudiants du Master Professionnel en Sécurité des Systèmes
Numériques de la Chaire UNESCO d'Afrique centrale, et aux praticiens qui
exercent ou exerceront comme auditeurs de la protection des données,
délégués à la protection des données, juristes conformité, RSSI ou
responsables de traitement. Il suppose acquis le droit du numérique (UE
MSN-2S3-JUR01) et la méthode générale de l'audit de cybersécurité (UE
MSN-2S3-SEC08). Il laisse au cours de forensique (UE MSN-2S3-SEC07)
l'enquête sur l'incident lui-même.

Trois partis pris le traversent. Le premier est de lire la loi dans son
texte officiel, article par article, sans lui prêter ce qu'elle ne dit pas~:
elle ne crée pas de délégué à la protection des données, elle ne fixe pas de
délai chiffré de notification, et de nombreuses dispositions attendent leurs
textes d'application. Le deuxième est de séparer, à chaque constat, le fait
établi de son interprétation, avec un niveau de confiance. Le troisième est
de coter ce que l'on ignore~: l'instrument \OUI{}~/ \NON{}~/ \NEK{} du
programme de recherche de l'auditologie sert ici à dire, pour chaque
exigence, si la conformité est établie, si l'écart est établi, ou si les
preuves manquent pour trancher. Une loi dont certains décrets ne sont pas
publiés produit des \NEK{} que l'auditeur doit assumer.

Les trois organisations qui accompagnent le lecteur, NYANGA CRÉDIT, le
CENTRE MÉDICAL MBOLO et NOUN PAY, sont fictives, comme leurs dirigeants,
leurs prestataires, les dates, les adresses et les montants. Les mécanismes,
les normes et les textes cités sont exacts. Les articles de loi sont cités
dans leur texte officiel~; les normes sous licence sont paraphrasées.

Le livre est accompagné d'un livret des corrigés et de NDAP SUKULU, où les
fiches de questions et les dix laboratoires du cours permettent de pratiquer
chaque chapitre.

\bigskip
\hfill\begin{minipage}{0.5\textwidth}
\raggedleft
\auteurcours\\
\laboratoire, Yaoundé, novembre 2026
\end{minipage}
